% 第10批：《大学物理简明教程（第5版）》第10章填空题10.2.1—10.2.20。
% 源扫描印刷/PDF第271—273页；题意简述、推导与交叉检查。

\subsection{填空题：光程、等厚干涉、夫琅禾费衍射和偏振}

\sourceproblem{10.2.1}{271--272}{271--272}
\problemtext 单色光从空气折射进入水中，频率、速度、波长各怎样改变？
\workings
介质界面处相位连续，光频率$f$保持不变。
水中相速度$v=c/n_{\rm water}$小于空气中速度，介质内波长$\lambda=v/f$也变短。
故填\textbf{不变、变小、变小}。
\crosscheck 不要把折射率增大误认为频率增大；光源频率决定时间周期。

\sourceproblem{10.2.2}{272}{272}
\problemtext 杨氏双缝波长$\lambda=562.5$ nm、缝屏距$D=1.2$ m，
相邻明纹间距$\Delta x=1.5$ mm，求缝间距。
\workings
由$\Delta x=\lambda D/d$，
\[
\boxed{d=\frac{562.5\times10^{-9}(1.2)}{1.5\times10^{-3}}
=4.50\times10^{-4}\,\mathrm m=0.450\,\mathrm{mm}}.
\]
\crosscheck 双缝间距大于光波长很多，旁轴条纹公式的几何光学假设合理。

\sourceproblem{10.2.3}{272}{272}
\problemtext 双缝形成明暗相间条纹，屏上$P$点位于第1级暗纹，求到该点的两束光程差。
\workings
暗纹满足$\Delta=(m+\tfrac12)\lambda$，以中心附近首条暗纹取$m=0$，
\[
\boxed{|\Delta|=\frac{\lambda}{2}}.
\]
\crosscheck 屏的两侧暗纹的光程差代数符号相反，但绝对值相同。

\sourceproblem{10.2.4}{272}{272}
\problemtext 用波长$600$ nm的单色光观察牛顿环，
求第2与第5条明环空气膜厚度之差。
\workings
反射明环两束光因一次半波损失，明纹所需膜厚呈公差$\lambda/2$的等差序列。
第5与第2明环相差三级，故
\[
\boxed{\Delta t=(5-2)\frac{\lambda}{2}=900\,\mathrm{nm}}.
\]
\crosscheck 求\emph{膜厚差}而非环半径差；公式与透镜曲率半径无关。

\sourceproblem{10.2.5}{272}{272}
\problemtext 白光照射牛顿环装置，考虑同一级条纹，问最靠近中心的是何种颜色。
\workings
对固定级次，反射环半径平方$ r_m^2\propto\lambda$；
短波长的环半径较小，所以最靠近中心的是\textbf{紫光}。
\crosscheck 同一可见级次中红光波长最长、环半径更大。

\sourceproblem{10.2.6}{272}{272}
\problemtext 杨氏双缝装置保持缝距与屏距不变，把全部装置由空气浸入水中，条纹变疏还是变密？
\workings
水中波长$\lambda'=\lambda/n_{\rm water}$，
$\Delta x'=\lambda'D/d=\Delta x/n_{\rm water}<\Delta x$，
故条纹\textbf{变密}。
\crosscheck 这里装置尺寸固定、光源频率固定；不能使用真空波长不变的错误假设。

\sourceproblem{10.2.7}{272}{272}
\problemtext 双缝干涉中只在下方狭缝覆盖一透明薄云母片，干涉条纹整体向哪里移动？
\workings
云母增加下方光路的光程$\delta=(n-1)t>0$。
为恢复两束光程相等，屏上零级明纹应移动到距离下方狭缝更近的位置，补偿几何光程差，
故条纹向\textbf{下方（被云母覆盖的狭缝方向）}移动。
\crosscheck 同时覆盖两条缝相同厚度不会造成这种整体平移。

\sourceproblem{10.2.8}{272}{272}
\problemtext 两块玻璃板构成左边接触的空气劈尖，
上玻璃板整体垂直于下玻璃板上移而不旋转，
反射条纹向哪边移动、间距怎样变？
\workings
若原膜厚$t(x)=\alpha x$，上板整体升高$\delta$后
$t'(x)=\alpha x+\delta$。
固定干涉级次要求$x_m'=(t_m-\delta)/\alpha<x_m$，
故纹向\textbf{棱边（较薄端）}平移。
斜率$\alpha$没有改变，条纹间距$\lambda/(2\alpha)$\textbf{不变}。
\crosscheck 若上板同时倾斜使$\alpha$改变则间距会变化，但不属于本题纯平移设定。

\sourceproblem{10.2.9}{272}{272}
\problemtext 单缝正入射，第一条暗纹的衍射角绝对值为$\varphi$，波长$\lambda$，求缝宽。
\workings
$a\sin\varphi=m\lambda$，首条暗纹$m=1$，故
\[
\boxed{a=\frac{\lambda}{\sin\varphi}}.
\]
\crosscheck 这是\emph{暗纹}的位置，不是一级明纹极大条件。

\sourceproblem{10.2.10}{272}{272}
\problemtext 单缝宽$a=4\lambda$，波长$\lambda$垂直入射，观察方向$\varphi=30^\circ$；波面分为几个半波带？
\workings
对应光程差$a\sin30^\circ=2\lambda$，
每个半波带对应$\lambda/2$，故
\[
\boxed{N=\frac{2a\sin\varphi}{\lambda}=4}.
\]
\crosscheck $4$个半波带可成对抵消，此方向同时是单缝第二级暗纹。

\sourceproblem{10.2.11}{272}{272}
\problemtext 单缝夫琅禾费衍射中减小缝宽、或增大入射光波长，各使条纹变疏还是变密？
\workings
暗纹角间隔尺度$\Delta\theta\sim\lambda/a$。
缝宽减小时增大，波长增大时也增大，
因此两空均填\textbf{变疏}。
\crosscheck 两项改变都增加中央明纹的角宽度。

\sourceproblem{10.2.12}{272}{272}
\problemtext 波长$600$ nm的单色光正入射空气中独立放置、
折射率$n_f=1.33$的透明薄膜；求反射光尽量相长的最小正膜厚。
若同样的膜改为\emph{镀在折射率$n_s=1.5$的玻璃透镜上}，
再求最小正膜厚。
\workings
(1) 对\emph{空气--膜--空气}的独立薄膜，
空气$\to$膜的第一反射发生一次$\pi$相变，
膜$\to$空气的第二反射不发生相变，
两束光之间有一次\emph{相对}半波损失。
反射光相长需
$2n_ft=(m+\tfrac12)\lambda$。
取$m=0$得
\[
\boxed{t_{1,\min}=\frac{\lambda}{4n_f}
=\frac{600}{4(1.33)}\,\mathrm{nm}\approx113\,\mathrm{nm}}.
\]
(2) 镀在玻璃透镜上时，
空气$\to$膜与膜$\to$玻璃都由低折射率向高折射率反射，
两束反射光均有$\pi$相变，\emph{相对}半波损失消失。
相长需$2n_ft=m\lambda$，
最小正$m=1$给
\[
\boxed{t_{2,\min}=\frac{\lambda}{2n_f}
=\frac{600}{2(1.33)}\,\mathrm{nm}\approx226\,\mathrm{nm}}.
\]
\crosscheck 两种薄膜的\emph{上表面材料相同}，差异在于下表面
是空气还是更高折射率的玻璃；正因这一区别，
最小反射增强膜厚相差两倍。

\sourceproblem{10.2.13}{272}{272}
\problemtext 单缝夫琅禾费衍射中，将缝沿凸透镜光轴移动，或沿垂直光轴方向平移，
焦平面上的衍射图样整体怎样移动？
\workings
在理想平面波入射、透镜焦面观察条件下，
远场强度分布是狭缝孔径函数傅里叶变换的模平方。
缝的平移引入复相位因子，但不改变其模平方。
因此两种平移在理想模型中均\textbf{不改变焦平面强度图样的中心位置}。
\crosscheck 若入射波不是准直平面波、或相对透镜改变有限距离使近轴成像条件失效，则需另行分析。

\sourceproblem{10.2.14}{272}{272}
\problemtext 单缝中央明纹用波长$589$ nm钠黄光时宽$4.0$ mm，
换成$442$ nm蓝紫光而装置不变，求中央明纹宽度。
\workings
中央明纹宽度与波长成正比，故
\[
\boxed{w_2=w_1\frac{\lambda_2}{\lambda_1}
=4.0\frac{442}{589}\approx3.00\,\mathrm{mm}}.
\]
\crosscheck 保持狭缝宽度与透镜焦距不变才能直接按波长比缩放。

\sourceproblem{10.2.15}{272}{272}
\problemtext 波长500 nm的平行光入射光栅，常量$d=2\times10^{-3}$ mm，
透光缝宽$a=1\times10^{-3}$ mm。问屏上可观察到多少条光栅主明纹。
\workings
$d=2\,\mu$m，$a=1\,\mu$m，$d/a=2$，所有偶数级$k=\pm2,\pm4,\ldots$与单缝零点重合而缺级。
由于$|k|\lambda/d<1$为可在有限焦面上观测的衍射角条件，
$|k|<4$，因此保留$k=0,\pm1,\pm3$共
\[
\boxed{N_{\mathrm{bright}}=5}.
\]
即屏上可观察到\textbf{5条主明纹}。
\]
\crosscheck $k=\pm4$仅在极限$\sin\theta=\pm1$的掠射方向，通常不计入焦平面屏幕可见条纹；
而且这里它们恰为缺级。

\sourceproblem{10.2.16}{272}{272}
\problemtext 马吕斯定律$I=I_0\cos^2\alpha$中，$I,I_0,\alpha$分别表示什么？
\workings
$I_0$为入射到检偏器的\emph{线偏振光强度}，
$I$为通过检偏器后的强度，
$\alpha$为入射线偏振光的振动方向与检偏器透振方向夹角。
\[
\boxed{I=I_0\cos^2\alpha}.
\]
\crosscheck 若$I_0$指自然光进入\emph{第一片}起偏器之前的强度，还需额外乘$1/2$。

\sourceproblem{10.2.17}{272}{272}
\problemtext 光强计中的检偏器旋转时，分别对自然光、线偏振光和部分偏振光判断光强变化的规律。
\workings
自然光经旋转检偏器的透射强度始终为入射的$1/2$，\textbf{不变}；
线偏振光满足马吕斯定律，最大最小交替，某些角度有\textbf{完全消光}；
部分偏振光既有自然光底噪又有线偏振成分，光强\textbf{周期变化但通常不完全消光}。
\crosscheck 只有完全线偏振光才能用理想检偏器出现严格零强度。

\sourceproblem{10.2.18}{272}{272}
\problemtext 自然光按布儒斯特角入射两介质分界面，反射光处于何种偏振状态，其振动方向与入射面何种关系？
\workings
布儒斯特反射光\textbf{完全线偏振}，电矢量的振动方向\textbf{垂直于入射面}。
\crosscheck 尽管折射光也发生部分偏振，但一般不是完全线偏振。

\sourceproblem{10.2.19}{272}{272}
\problemtext 自然光从空气入射平板玻璃，反射光完全偏振且\emph{折射角}$r=30^\circ$，空气折射率取1，求玻璃折射率。
\workings
布儒斯特条件$i_B+r=90^\circ$，故$i_B=60^\circ$，
并有$n_2/n_1=\tan i_B$，所以
\[
\boxed{n_{\mathrm{glass}}=\tan60^\circ=\sqrt3\approx1.732}.
\]
\crosscheck 原题给的是\emph{折射角}30°而不是\emph{入射角}30°；
互换会错误得出小于1的“玻璃折射率”。

\sourceproblem{10.2.20}{273}{273}
\problemtext 光的干涉和衍射说明光具有何种性质？偏振现象说明光属于哪种波？
\workings
干涉和衍射是\textbf{波动性}的证据；
偏振表明电场振动方向与光传播方向垂直，因此光是\textbf{横波}。
\crosscheck 波动性和横波属性是两项不同层次的结论，不能把两空都填“电磁波”而丢失题意。
